\documentclass[11pt]{article}
\usepackage[margin=25mm]{geometry}
\usepackage{fontspec}
\setmainfont{Latin Modern Roman}
\usepackage{amsmath,amssymb,mathtools}
\usepackage{unicode-math}
\setmathfont{Latin Modern Math}
\usepackage{xurl,hyperref}
\hypersetup{colorlinks=true,urlcolor=blue}
\setcounter{secnumdepth}{0}
\setlength{\parindent}{0pt}
\setlength{\parskip}{5pt}
\setlength{\emergencystretch}{3em}
\newcommand{\sref}[2]{\hyperlink{source:#1:#2}{[S#2]}}
\newcommand{\cataloguescope}[1]{\paragraph{Recorded target.}#1}
\begin{document}
\section{No 5 Points on 2-plane in $[N]^d$}\label{statement}

\subsection{Definitions and mathematical statement}

\paragraph{Shared notation.}$\mathbb N=\{1,2,\ldots\}$, and $\mathbb Z,\mathbb Q,\mathbb R,\mathbb C$ denote the integers, rationals, reals and complex numbers. Any other conventions are specified locally.


For integers $N\geq1$ and $d\geq2$, put $[N]=\{1,\ldots,N\}$ and view $[N]^d$ as a subset of $\mathbb R^d$. Let $\mathcal P_2(\mathbb R^d)$ denote the affine two-dimensional planes, and define
\[
 M_d(N)=\max\bigl\{|A|:A\subset[N]^d,\ |A\cap P|\leq4\text{ for every }P\in\mathcal P_2(\mathbb R^d)\bigr\}.
\]
The condition includes five collinear points, since a line is contained in an affine two-plane; it does not require the five points to span that plane.
\paragraph{Statement recorded in the supplied source.}
Determine the largest possible size $M_d(N)$, and obtain sharp estimates for its growth with $N$ and $d$. \sref{statement}{2}
\subsection{Short English statement}
How large can a subset of a high-dimensional integer grid be if every affine two-plane contains at most four of its points?
\subsection{Sources}
\begin{description}
\item[\hypertarget{source:statement:1}{S1}] ULAM.AI and the UnsolvedMath Contributors, \emph{UnsolvedMath: A Curated Collection of Open Mathematics Problems}, catalogue entry. CC BY 4.0. \url{https://huggingface.co/datasets/ulamai/UnsolvedMath}.
\item[\hypertarget{source:statement:2}{S2}] Ben Green, \emph{100 Open Problems}, maintained author collection, Problem 74 and its comments, accessed 7 October 2026. \url{https://people.maths.ox.ac.uk/greenbj/papers/open-problems.pdf}.
\end{description}
\label{end:statement}
\end{document}
